\ifdefined\gkver\else\def\gkver{student}\fi
\documentclass[\gkver]{gaokaozhenti}
\begin{document}

\chapter{1995年上海}

\begin{enumerate}
\item
%% number: 1
%% paperName: 1995年上海高考第 1 题 4 分
%% typeId: 1
%% type: 单选题
%% chapter: 光学
%% point: 光电效应
%% method: 无
%% score: 4
%% degree: 850
%% duplicateId: 0
%% body:
对光电效应作出合理解释的物理学家是\xzanswer{A}
\fourchoices[answer=A]
{爱因斯坦}
{玻尔}
{查德威克}
{德布罗意}
%% answer: A
%% memo:
\memoanswer{
本题原卷及材料中未提供详解，待补充。
}

\item
%% number: 2
%% paperName: 1995年上海高考第 2 题 4 分
%% typeId: 1
%% type: 单选题
%% chapter: 交流电
%% point: LC振荡回路
%% method: 无
%% score: 4
%% degree: 650
%% duplicateId: 0
%% body:
在 $LC$ 振荡电路中，电容器上的带电量从最大值变化到零所需的最短时间是\xzanswer{B}
\fourchoices[answer=B]
{$\frac{1}{4}\pi\sqrt{LC}$}
{$\frac{1}{2}\pi\sqrt{LC}$}
{$\pi\sqrt{LC}$}
{$2\pi\sqrt{LC}$}
%% answer: B
%% memo:
\memoanswer{
本题原卷及材料中未提供详解，待补充。
}

\item
%% number: 3
%% paperName: 1995年上海高考第 3 题 4 分
%% typeId: 1
%% type: 单选题
%% chapter: 气体性质
%% point: 热力学第一定律
%% method: 无
%% score: 4
%% degree: 650
%% duplicateId: 0
%% body:
关于物体的内能变化，以下说法中正确的是\xzanswer{C}
\fourchoices[answer=C]
{物体吸收热量，内能一定增大}
{物体对外做功，内能一定减小}
{物体吸收热量，同时对外做功，内能可能不变}
{物体放出热量，同时对外做功，内能可能不变}
%% answer: C
%% memo:
\memoanswer{
本题原卷及材料中未提供详解，待补充。
}

\item
%% number: 4
%% paperName: 1995年上海高考第 4 题 4 分
%% typeId: 1
%% type: 单选题
%% chapter: 电路
%% point: 电功电功率
%% method: 对称法
%% score: 4
%% degree: 650
%% duplicateId: 0
%% body:
如图所示电路中，六个电阻的阻值均相同，由于对称性，电阻 $R_{2}$ 上无电流流过。已知电阻 $R_{6}$ 所消耗的电功率为 $1\UW$，则六个电阻所消耗的总功率为\xzanswer{D}
\onepicture[w=6cm]{\includesvg{figs/04}}
\fourchoices[answer=D]
{$6\UW$}
{$5\UW$}
{$3\UW$}
{$2\UW$}
%% answer: D
%% memo:
\memoanswer{
本题原卷及材料中未提供详解，待补充。
}

\item
%% number: 5
%% paperName: 1995年上海高考第 5 题 4 分
%% typeId: 1
%% type: 单选题
%% chapter: 功和能
%% point: 动能定理
%% method: 无
%% score: 4
%% degree: 450
%% duplicateId: 0
%% body:
一物体静止在光滑水平面上，先对物体施一水平向右的恒力 $F_{1}$，经 $t$ 秒后撤去 $F_{1}$，立即再对它施一水平向左的恒力 $F_{2}$，又经 $t$ 秒后物体回到出发点。在这一过程中，$F_{1}$、$F_{2}$ 分别对物体作的功 $W_{1}$、$W_{2}$ 间的关系是\xzanswer{C}
\fourchoices[answer=C]
{$W_{2} = W_{1}$}
{$W_{2} = 2W_{1}$}
{$W_{2} = 3W_{1}$}
{$W_{2} = 5W_{1}$}
%% answer: C
%% memo:
\memoanswer{
本题原卷及材料中未提供详解，待补充。
}

\item
%% number: 6
%% paperName: 1995年上海高考第 6 题 4 分
%% typeId: 1
%% type: 单选题
%% chapter: 气体性质
%% point: 气体的状态参量
%% method: 微元
%% score: 4
%% degree: 450
%% duplicateId: 0
%% body:
两个半球壳拼成的球形容器内部已抽成真空，球形容器的半径为 $R$，大气压强为 $p$。为使两个半球壳沿图中箭头方向互相分离，应施加的力 $F$ 至少为\xzanswer{C}
\onepicture[w=7cm]{\includesvg{figs/06}}
\fourchoices[answer=C]
{$4\pi R^{2}p$}
{$2\pi R^{2}p$}
{$\pi R^{2}p$}
{$\frac{\pi R^{2}p}{2}$}
%% answer: C
%% memo:
\memoanswer{
本题原卷及材料中未提供详解，待补充。
}

\item
%% number: 7
%% paperName: 1995年上海高考第 7 题 4 分
%% typeId: 1
%% type: 单选题
%% chapter: 力物体的平衡
%% point: 三力平衡
%% method: 无
%% score: 4
%% degree: 650
%% duplicateId: 0
%% body:
三个相同的支座上分别搁着三个质量和直径都相等的光滑圆球 a、b、c，支点 P、Q 在同一水平面上，a 球的重心 $O_{a}$ 位于球心，b 球和 c 球的重心 $O_{b}$、$O_{c}$ 分别位于球心的正上方和球心的正下方，如图所示。三球均处于平衡状态。支点 P 对 a 球的弹力为 $N_{a}$，对 b 球和 c 球的弹力分别为 $N_{b}$ 和 $N_{c}$，则\xzanswer{A}
\onepicture[w=10cm]{\includesvg{figs/07}}
\fourchoices[answer=A]
{$N_{a} = N_{b} = N_{c}$}
{$N_{a} > N_{b} > N_{c}$}
{$N_{a} < N_{b} < N_{c}$}
{$N_{a} > N_{b} = N_{c}$}
%% answer: A
%% memo:
\memoanswer{
本题原卷及材料中未提供详解，待补充。
}

\item
%% number: 8
%% paperName: 1995年上海高考第 8 题 4 分
%% typeId: 1
%% type: 单选题
%% chapter: 气体性质
%% point: 变速运动中的压强
%% method: 无
%% score: 4
%% degree: 450
%% duplicateId: 0
%% body:
如图所示，a，b，c 三根完全相同的玻璃管，一端封闭，管内各用相同长度的一段水银柱封闭了质量相等的空气。a 管竖直向下作自由落体运动，b 管竖直向上作加速度为 $g$ 的匀加速运动，c 管沿倾角为 $45^{\circ}$ 的光滑斜面下滑。若空气温度始终不变，当水银柱相对管壁静止时，a，b，c 三管内的空气柱长度 $l_{a}$、$l_{b}$、$l_{c}$ 间的关系为\xzanswer{D}
\onepicture[w=7cm]{\includesvg{figs/08}}
\fourchoices[answer=D]
{$l_{b} = l_{c} = l_{a}$}
{$l_{b} < l_{c} < l_{a}$}
{$l_{b} > l_{c} > l_{a}$}
{$l_{b} < l_{c} = l_{a}$}
%% answer: D
%% memo:
\memoanswer{
本题原卷及材料中未提供详解，待补充。
}

\item
%% number: 9
%% paperName: 1995年上海高考第 9 题 5 分
%% typeId: 2
%% type: 多选题
%% chapter: 气体性质
%% point: 理想气体状态方程
%% method: 无
%% score: 5
%% degree: 650
%% duplicateId: 0
%% body:
一定质量的理想气体发生状态变化时，其状态量 $p$、$V$、$T$ 的变化情况可能是\xzanswer{ABD}
\fourchoices[answer=ABD]
{$p$、$V$、$T$ 都增大}
{$p$ 减小，$V$ 和 $T$ 增大}
{$p$ 和 $V$ 减小，$T$ 增大}
{$p$ 和 $T$ 增大，$V$ 减小}
%% answer: ABD
%% memo:
\memoanswer{
本题原卷及材料中未提供详解，待补充。
}

\item
%% number: 10
%% paperName: 1995年上海高考第 10 题 5 分
%% typeId: 2
%% type: 多选题
%% chapter: 交流电
%% point: 变压器
%% method: 无
%% score: 5
%% degree: 650
%% duplicateId: 0
%% body:
如图所示，理想变压器的副线圈上通过输电线接有两个相同的灯泡 $L_{1}$ 和 $L_{2}$，输电线的等效电阻为 $R$。开始时，电键 K 断开。当 K 接通时，以下说法中正确的是\xzanswer{BCD}
\onepicture[w=8cm]{\includesvg{figs/10}}
\fourchoices[answer=BCD]
{副线圈两端 M、N 的输出电压减小}
{副线圈输电线等效电阻 $R$ 上的电压降增大}
{通过灯泡 $L_{1}$ 的电流减小}
{原线圈中的电流增大}
%% answer: BCD
%% memo:
\memoanswer{
本题原卷及材料中未提供详解，待补充。
}

\item
%% number: 11
%% paperName: 1995年上海高考第 11 题 5 分
%% typeId: 2
%% type: 多选题
%% chapter: 直线运动
%% point: 时间中点速度
%% method: 图象法
%% score: 5
%% degree: 650
%% duplicateId: 0
%% body:
物体沿一直线运动，在 $t$ 时间内通过的路程为 $s$。它在中间位置 $s/2$ 处的速度为 $v_{1}$，在中间时刻 $t/2$ 时的速度为 $v_{2}$，则 $v_{1}$ 和 $v_{2}$ 的关系为\xzanswer{ABC}
\fourchoices[answer=ABC]
{当物体作匀加速直线运动时，$v_{1} > v_{2}$}
{当物体作匀减速直线运动时，$v_{1} > v_{2}$}
{当物体作匀速直线运动时，$v_{1} = v_{2}$}
{当物体作匀减速直线运动时，$v_{1} < v_{2}$}
%% answer: ABC
%% memo:
\memoanswer{
本题原卷及材料中未提供详解，待补充。
}

\item
%% number: 12
%% paperName: 1995年上海高考第 12 题 5 分
%% typeId: 2
%% type: 多选题
%% chapter: 电场
%% point: 带电粒子的平衡
%% method: 无
%% score: 5
%% degree: 650
%% duplicateId: 0
%% body:
如图所示，两板间距为 $d$ 的平行板电容器与一电源连接，电键 K 闭合。电容器两板间有一质量为 $m$、带电量为 $q$ 的微粒静止不动。下列各叙述中正确的是\xzanswer{BD}
\onepicture[w=7cm]{\includesvg{figs/12}}
\fourchoices[answer=BD]
{微粒带的是正电}
{电源电动势的大小等于 $\frac{mgd}{q}$}
{断开电键 K，微粒将向下作加速运动}
{保持电键 K 闭合，把电容器两极板距离增大，微粒将向下作加速运动}
%% answer: BD
%% memo:
\memoanswer{
本题原卷及材料中未提供详解，待补充。
}

\item
%% number: 13
%% paperName: 1995年上海高考第 13 题 5 分
%% typeId: 2
%% type: 多选题
%% chapter: 电磁感应
%% point: 电磁感应受力分析
%% method: 无
%% score: 5
%% degree: 450
%% duplicateId: 0
%% body:
如图所示，通有稳恒电流的螺线管竖直放置，铜环 $R$ 沿螺线管的轴线加速下落。在下落过程中，环面始终保持水平。铜环先后经过轴线上 1、2、3 位置时的加速度分别为 $a_{1}$、$a_{2}$、$a_{3}$。位置 2 处于螺线管的中心，位置 1、3 与位置 2 等距离。则\xzanswer{ABD}
\onepicture[h=5cm]{\includesvg{figs/13}}
\fourchoices[answer=ABD]
{$a_{1} < a_{2} = g$}
{$a_{3} < a_{1} < g$}
{$a_{1} = a_{3} < a_{2}$}
{$a_{3} < a_{1} < a_{2}$}
%% answer: ABD
%% memo:
\memoanswer{
本题原卷及材料中未提供详解，待补充。
}

\item
%% number: 14
%% paperName: 1995年上海高考第 14 题 4 分
%% typeId: 3
%% type: 填空题
%% chapter: 电路
%% point: 等效电路
%% method: 无
%% score: 4
%% degree: 650
%% duplicateId: 0
%% body:
（A 组）如图所示电路中，电阻 $R_{1}$、$R_{2}$、$R_{3}$ 的阻值都是 $1\UO$，$R_{4}$、$R_{5}$ 的阻值都是 $0.5\UO$，ab 端输入电压 $U = 6\UV$。当 cd 端接伏特计时，其示数是 \tkanswer[1.6cm]{$2\UV$}。
\onepicture[align=right, w=6cm]{\includesvg{figs/14}}
%% answer: 
\jdanswer{
$2\UV$
}
%% memo:
\memoanswer{
本题原卷及材料中未提供详解，待补充。
}

\item
%% number: 14
%% paperName: 1995年上海高考第 14 题 4 分
%% typeId: 3
%% type: 填空题
%% chapter: 电路
%% point: 等效电路
%% method: 无
%% score: 4
%% degree: 650
%% duplicateId: 0
%% body:
（B 组）如图所示电路中，电阻 $R_{1}$、$R_{2}$、$R_{3}$ 的阻值都是 $1\UO$，$R_{4}$、$R_{5}$ 的阻值都是 $0.5\UO$，ab 端输入电压 $U = 5\UV$。当 cd 端接安培计时，其示数是 \tkanswer[1.6cm]{$1\UA$}。
\onepicture[align=right, w=6cm]{\includesvg{figs/14}}
%% answer: 
\jdanswer{
$1\UA$
}
%% memo:
\memoanswer{
本题原卷及材料中未提供详解，待补充。
}

\item
%% number: 15
%% paperName: 1995年上海高考第 15 题 4 分
%% typeId: 3
%% type: 填空题
%% chapter: 原子和原子核
%% point: 质能方程
%% method: 无
%% score: 4
%% degree: 650
%% duplicateId: 0
%% body:
（A 组）质子、中子和氘核的质量分别为 $m_{1}$、$m_{2}$ 和 $m_{3}$。质子和中子结合成氘核时，发出 $\gamma$ 射线。已知普朗克恒量为 $h$，真空中光速为 $c$，则 $\gamma$ 射线的频率 $\nu =$ \tkanswer[3.2cm]{$\frac{(m_{1} + m_{2} - m_{3})c^{2}}{h}$}。
%% answer: 
\jdanswer{
$\frac{(m_{1} + m_{2} - m_{3})c^{2}}{h}$
}
%% memo:
\memoanswer{
本题原卷及材料中未提供详解，待补充。
}

\item
%% number: 15
%% paperName: 1995年上海高考第 15 题 4 分
%% typeId: 3
%% type: 填空题
%% chapter: 原子和原子核
%% point: 玻尔模型
%% method: 无
%% score: 4
%% degree: 650
%% duplicateId: 0
%% body:
（B 组）按照玻尔理论，氢原子处在量子数为 $n = 2$ 和 $n = 3$ 的定态时，其相应的原子能量的绝对值之比 $|E_{2}|:|E_{3}| =$ \tkanswer[2cm]{$9:4$}。
%% answer: 
\jdanswer{
$9:4$
}
%% memo:
\memoanswer{
本题原卷及材料中未提供详解，待补充。
}

\item
%% number: 16
%% paperName: 1995年上海高考第 16 题 4 分
%% typeId: 3
%% type: 填空题
%% chapter: 光学
%% point: 光的折射
%% method: 无
%% score: 4
%% degree: 650
%% duplicateId: 0
%% body:
（A 组）有一半圆形玻璃砖，玻璃折射率为 $\sqrt{2}$，AB 为其直径，O 为圆心。一束宽度恰等于玻璃砖半径的单色平行光垂直于 AB 从空气射入玻璃砖，其中心光线通过 O 点。则光束中的光线射出玻璃砖时最大的折射角为 \tkanswer[1.6cm]{$45^{\circ}$}。并画出图中三条光线在玻璃砖内和射出玻璃砖后的光路。
\onepicture[align=right, h=4.2cm]{\includesvg{figs/16a}}
%% answer: 
\jdanswer{
\begin{enumerate}
	\item
	$45^{\circ}$

	\item
	\onepicture[w=6cm]{\includesvg{figs/16a_an}}

\end{enumerate}
}
%% memo:
\memoanswer{
本题原卷及材料中未提供详解，待补充。
}

\item
%% number: 16
%% paperName: 1995年上海高考第 16 题 4 分
%% typeId: 3
%% type: 填空题
%% chapter: 光学
%% point: 透镜
%% method: 无
%% score: 4
%% degree: 650
%% duplicateId: 0
%% body:
（B 组）在焦距为 $f$ 的凸透镜的主光轴上，与光心 O 相距 $2f$ 处放一点光源 S。已知凸透镜的半径等于 $R$。若在凸透镜另一侧的焦点 F 处垂直于主光轴放一光屏，则在光屏上呈现的圆形光斑的面积为 \tkanswer[2.4cm]{$\frac{1}{4}\pi R^{2}$}。并画出图中自 S 发出的三条光线经凸透镜后的光路。
\onepicture[align=right, w=6.5cm]{\includesvg{figs/16b}}
%% answer: 
\jdanswer{
\begin{enumerate}
	\item
	$\frac{1}{4}\pi R^{2}$

	\item
	\onepicture[w=6.5cm]{\includesvg{figs/16b_an}}

\end{enumerate}
}
%% memo:
\memoanswer{
本题原卷及材料中未提供详解，待补充。
}

\item
%% number: 17
%% paperName: 1995年上海高考第 17 题 4 分
%% typeId: 3
%% type: 填空题
%% chapter: 气体性质
%% point: 阿伏伽德罗常数
%% method: 无
%% score: 4
%% degree: 650
%% duplicateId: 0
%% body:
水分子的质量等于 \tkanswer[3cm]{$2.99\times10^{-26}\Ukg$}。已知阿伏伽德罗常数为 $6.02\times10^{23}\Umoln$。
%% answer: 
\jdanswer{
$2.99\times10^{-26}\Ukg$
}
%% memo:
\memoanswer{
本题原卷及材料中未提供详解，待补充。
}

\item
%% number: 18
%% paperName: 1995年上海高考第 18 题 4 分
%% typeId: 3
%% type: 填空题
%% chapter: 运动和力
%% point: 牛顿第二定律
%% method: 临界
%% score: 4
%% degree: 450
%% duplicateId: 0
%% body:
如图所示，一细线的一端固定于倾角为 $45^{\circ}$ 的光滑楔形滑块 A 的顶端 P 处，细线的另一端拴一质量为 $m$ 的小球。当滑块至少以加速度 $a =$ \tkanswer[1.4cm]{$g$} 向左运动时，小球对滑块的压力等于零。当滑块以 $a = 2g$ 的加速度向左运动时，线中拉力 $T =$ \tkanswer[2cm]{$\sqrt{5}mg$}。
\onepicture[align=right, w=5.5cm]{\includesvg{figs/18}}
%% answer: 
\jdanswer{
$g$，$\sqrt{5}mg$
}
%% memo:
\memoanswer{
本题原卷及材料中未提供详解，待补充。
}

\item
%% number: 19
%% paperName: 1995年上海高考第 19 题 4 分
%% typeId: 3
%% type: 填空题
%% chapter: 电场
%% point: 带电粒子的圆周运动
%% method: 等效替代
%% score: 4
%% degree: 650
%% duplicateId: 0
%% body:
如图所示，一绝缘细圆环半径为 $r$，其环面固定在水平面上。场强为 $E$ 的匀强电场与圆环平面平行。环上穿有一电量为 $+q$、质量为 $m$ 的小球，可沿圆环做无摩擦的圆周运动。若小球经 A 点时速度 $v_{A}$ 的方向恰与电场垂直，且圆环与小球间沿水平方向无力的作用，则速度 $v_{A} =$ \tkanswer[2.4cm]{$\sqrt{\frac{qEr}{m}}$}。当小球运动到与 A 点对称的 B 点时，小球对圆环在水平方向的作用力 $N_{B} =$ \tkanswer[1.6cm]{$6qE$}。
\onepicture[align=right, w=5cm]{\includesvg{figs/19}}
%% answer: 
\jdanswer{
$\sqrt{\frac{qEr}{m}}$，$6qE$
}
%% memo:
\memoanswer{
本题原卷及材料中未提供详解，待补充。
}

\item
%% number: 20
%% paperName: 1995年上海高考第 20 题 4 分
%% typeId: 3
%% type: 填空题
%% chapter: 电场
%% point: 电场叠加
%% method: 对称法
%% score: 4
%% degree: 450
%% duplicateId: 0
%% body:
一半径为 $R$ 的绝缘球壳上均匀地带有电量为 $+Q$ 的电荷，另一电量为 $+q$ 的点电荷放在球心 O 上，由于对称性，点电荷受力为零。现在球壳上挖去半径为 $r$（$r$ 远小于 $R$）的一个小圆孔，则此时置于球心的点电荷所受力的大小为 \tkanswer[3cm]{$\frac{kQq r^{2}}{4R^{2}}$}（已知静电力恒量为 $k$），方向 \tkanswer[2.6cm]{由球心指向小孔中心}。
%% answer: 
\jdanswer{
$\frac{kQq r^{2}}{4R^{2}}$，由球心指向小孔中心
}
%% memo:
\memoanswer{
本题原卷及材料中未提供详解，待补充。
}

\item
%% number: 21
%% paperName: 1995年上海高考第 21 题 4 分
%% typeId: 3
%% type: 填空题
%% chapter: 机械波
%% point: 波的叠加
%% method: 无
%% score: 4
%% degree: 450
%% duplicateId: 0
%% body:
两列简谐横波均沿 $x$ 轴传播，传播速度的大小相等，其中一列沿正 $x$ 方向传播（图中实线所示），一列沿负 $x$ 方向传播（图中虚线所示）。这两列波的频率相等，振动方向均沿 $y$ 轴，则图中 $x = 1, 2, 3, 4, 5, 6, 7, 8$ 各点中振幅最大的是 $x =$ \tkanswer[2cm]{$4$ 和 $8$} 的点，振幅最小的是 $x =$ \tkanswer[2cm]{$2$ 和 $6$} 的点。
\onepicture[align=right, w=12cm]{\includesvg{figs/21}}
%% answer: 
\jdanswer{
$4$ 和 $8$，$2$ 和 $6$
}
%% memo:
\memoanswer{
本题原卷及材料中未提供详解，待补充。
}

\item
%% number: 22
%% paperName: 1995年上海高考第 22 题 4 分
%% typeId: 4
%% type: 实验题
%% chapter: 磁场
%% point: 左手定则
%% method: 无
%% score: 4
%% degree: 650
%% duplicateId: 0
%% body:
如图所示，放在马蹄形磁铁两极之间的导体棒 ab，当通有自 b 到 a 的电流时受到向右的安培力作用，则磁铁的上端是 \tkanswer[1.4cm]{$N$} 极。如果磁铁上端是 S 极，导体棒中的电流方向自 a 到 b，则导体棒受到的安培力方向向 \tkanswer[1.4cm]{右}。
\onepicture[align=right, w=5cm]{\includesvg{figs/22}}
%% answer: 
\jdanswer{
$N$，右
}
%% memo:
\memoanswer{
本题原卷及材料中未提供详解，待补充。
}

\item
%% number: 23
%% paperName: 1995年上海高考第 23 题 5 分
%% typeId: 4
%% type: 实验题
%% chapter: 气体性质
%% point: 验证玻意耳定律
%% method: 无
%% score: 5
%% degree: 650
%% duplicateId: 0
%% body:
在“验证玻意耳定律”的实验中，对气体的初状态和末状态的测量和计算都正确无误。结果末状态的 $pV$ 值与初状态的 $p_{0}V_{0}$ 值明显不等。造成这一结果的可能原因是在实验过程中\xzanswer{ACD}
\fourchoices[answer=ACD]
{气体温度发生变化}
{气体与外界间有热交换}
{有气体泄漏}
{体积改变得太迅速}
%% answer: ACD
\jdanswer{
ACD
}
%% memo:
\memoanswer{
本题原卷及材料中未提供详解，待补充。
}

\item
%% number: 24
%% paperName: 1995年上海高考第 24 题 7 分
%% typeId: 4
%% type: 实验题
%% chapter: 电路
%% point: 测电动势与内阻
%% method: 无
%% score: 7
%% degree: 650
%% duplicateId: 0
%% body:
某学生用安培计和伏特计测干电池的电动势和内阻时，所用滑动变阻器的阻值范围为 $0\sim20\UO$，连接电路的实物图如下。
\onepicture[align=right, w=8cm]{\includesvg{figs/24}}
①该学生接线中错误的和不规范的做法是\xzanswer{ABD}
\fourchoices[answer=ABD]
{滑动变阻器不起变阻作用}
{安培计接线有错}
{伏特计量程选用不当}
{伏特计接线有错}
②在方框里画出这个实验的正确电路图。
%% answer: 
\jdanswer{
\begin{enumerate}
	\item
	①ABD

	\item
	②见图
	\onepicture[w=4.5cm]{\includesvg{figs/24_an}}

\end{enumerate}
}
%% memo:
\memoanswer{
本题原卷及材料中未提供详解，待补充。
}

\item
%% number: 25
%% paperName: 1995年上海高考第 25 题 7 分
%% typeId: 4
%% type: 实验题
%% chapter: 曲线运动
%% point: 研究平抛运动实验
%% method: 无
%% score: 7
%% degree: 450
%% duplicateId: 0
%% body:
试根据平抛运动原理设计测量弹射器弹丸出射初速的实验方法。提供实验器材：弹射器（含弹丸，见示意图）；铁架台（带有夹具）；米尺。
\onepicture[align=right, w=4cm]{\includesvg{figs/25}}
①在方框中画出实验示意图。

②在安装弹射器时应注意：\tkanswer[2.6cm]{弹射器必需保持水平}；

③实验中需要测量的量（并在示意图中用字母标出）：\tkanswer[3.4cm]{弹丸下降高度 $y$ 和水平射程 $x$}；

④由于弹射器每次射出的弹丸初速不可能完全相等，在实验中应采取的方法是：\tkanswer[4.4cm]{在不改变高度 $y$ 的条件下进行多次实验，测量水平射程 $x$，得出平均水平射程 $\bar{x}$}；

⑤计算公式：\tkanswer[2.8cm]{$\overline{v_{0}} = \frac{\bar{x}}{\sqrt{\frac{2y}{g}}}$}。
%% answer: 
\jdanswer{
\begin{enumerate}
	\item
	①见图
	\onepicture[w=7cm]{\includesvg{figs/25_an}}

	\item
	②弹射器必需保持水平。

	\item
	③弹丸下降高度 $y$ 和水平射程 $x$。

	\item
	④在不改变高度 $y$ 的条件下进行多次实验，测量水平射程 $x$，得出平均水平射程 $\bar{x}$。

	\item
	⑤$\overline{v_{0}} = \frac{\bar{x}}{\sqrt{\frac{2y}{g}}}$

\end{enumerate}
}
%% memo:
\memoanswer{
本题原卷及材料中未提供详解，待补充。
}

\item
%% number: 26
%% paperName: 1995年上海高考第 26 题 11 分
%% typeId: 6
%% type: 计算题
%% chapter: 电磁感应
%% point: 电磁感应电路分析
%% method: 无
%% score: 11
%% degree: 650
%% duplicateId: 0
%% body:
把总电阻为 $2R$ 的均匀电阻丝焊接成一半径为 $a$ 的圆环，水平固定在竖直向下的磁感应强度为 $B$ 的匀强磁场中，如图所示。一长度为 $2a$，电阻等于 $R$，粗细均匀的金属棒 MN 放在圆环上，它与圆环始终保持良好的电接触。当金属棒以恒定速度 $v$ 向右移动，经过环心 O 时，求：
\begin{enumerate}
	\item
	棒上电流的大小和方向，及棒两端的电压 $U_{MN}$；
	\item
	在圆环和金属棒上消耗的总热功率。
\end{enumerate}
\onepicture[align=right, h=5cm]{\includesvg{figs/26}}
%% answer: 
\jdanswer{
\begin{enumerate}
	\item
	$I = \frac{4Bav}{3R}$，方向由 N 到 M，$U_{MN} = \frac{2}{3}Bav$

	\item
	$P = \frac{8B^{2}a^{2}v^{2}}{3R}$

\end{enumerate}
}
%% memo:
\memoanswer{
（1）$E = 2Bav$，$I = \frac{E}{R + \frac{R}{2}} = \frac{2E}{3R} = \frac{4Bav}{3R}$，方向由 N 到 M。

$U_{MN} = I\cdot\frac{R}{2} = \frac{4Bav}{3R}\cdot\frac{R}{2} = \frac{2}{3}Bav$

（2）$P = IE = \frac{4Bav}{3R}\cdot 2Bav = \frac{8B^{2}a^{2}v^{2}}{3R}$

评分标准：（1）7 分。求出 $E$ 2 分。求出 $I$ 大小 2 分，正确指出方向 1 分，正确得出 $U_{MN}$ 2 分。

（2）4 分。若分别求出环上热功率 $P_{\text{环}}$ 和棒上热功率 $P_{\text{棒}}$。求对一个给 1 分，求出总和再给 2 分。
}

\item
%% number: 27
%% paperName: 1995年上海高考第 27 题 12 分
%% typeId: 6
%% type: 计算题
%% chapter: 光学
%% point: 光的折射
%% method: 无
%% score: 12
%% degree: 650
%% duplicateId: 0
%% body:
宽为 $a$ 的平行光束从空气斜向入射到两面平行的玻璃板上表面，入射角为 $45^{\circ}$，光束中包含两种波长的光，玻璃对这两种波长光的折射率分别为 $n_{1} = 1.5$，$n_{2} = \sqrt{3}$。
\begin{enumerate}
	\item
	求每种波长的光射入上表面后的折射角 $r_{1}$、$r_{2}$；
	\item
	为使光束从玻璃板下表面出射时能分成不交叠的两束，玻璃板的厚度 $d$ 至少为多少？并画出光路示意图。
\end{enumerate}
\onepicture[align=right, w=6.5cm]{\includesvg{figs/27}}
%% answer: 
\jdanswer{
\begin{enumerate}
	\item
	$r_{1} = \sin^{-1}\frac{\sqrt{2}}{3}$，$r_{2} = \sin^{-1}\frac{\sqrt{6}}{6}$

	\item
	$d = \frac{7\sqrt{10} + 10\sqrt{7}}{3}a$

\end{enumerate}
}
%% memo:
\memoanswer{
（1）由折射定律

$\sin r_{1} = \frac{\sin i}{n_{1}} = \frac{\frac{\sqrt{2}}{2}}{1.5} = \frac{\sqrt{2}}{3}$，解得 $r_{1} = \sin^{-1}\frac{\sqrt{2}}{3}$

$\sin r_{2} = \frac{\sin i}{n_{2}} = \frac{\frac{\sqrt{2}}{2}}{\sqrt{3}} = \frac{\sqrt{6}}{6}$，解得 $r_{2} = \sin^{-1}\frac{\sqrt{6}}{6}$

（2）$d\tan r_{1} - d\tan r_{2} = \frac{a}{\cos i}$     ①

\onepicture[align=right, w=6cm]{\includesvg{figs/27_an}}

$d = \frac{a}{\cos i(\tan r_{1} - \tan r_{2})}$     ②

$\tan r_{1} = \tan\left(\sin^{-1}\frac{\sqrt{2}}{3}\right) = \sqrt{\frac{2}{7}}$，$\tan r_{2} = \tan\left(\sin^{-1}\frac{\sqrt{6}}{6}\right) = \sqrt{\frac{1}{5}}$

代入②式解得：

$d = \frac{a}{\frac{\sqrt{2}}{2}\left(\sqrt{\frac{2}{7}} - \sqrt{\frac{1}{5}}\right)} = \frac{a}{\sqrt{\frac{1}{7}} - \sqrt{\frac{1}{10}}} = \frac{7\sqrt{10} + 10\sqrt{7}}{3}a$     ③

评分标准：（1）4 分。正确求出 $r_{1}$、$r_{2}$ 各 2 分。

（2）8 分。画出光路示意图得 2 分，正确列出①式（或②式）得 3 分。正确得出结果（③式）得 3 分。
}

\item
%% number: 28
%% paperName: 1995年上海高考第 28 题 15 分
%% typeId: 6
%% type: 计算题
%% chapter: 运动和力
%% point: 动量守恒定律
%% method: 无
%% score: 15
%% degree: 450
%% duplicateId: 0
%% body:
如图所示，A、B、C 三物块质量均为 $m$，置于光滑水平台面上。B、C 间夹有原已完全压紧不能再压缩的弹簧，两物块用细绳相连，使弹簧不能伸展。物块 A 以初速 $v_{0}$ 沿 B、C 连线方向向 B 运动，相碰后，A 与 B、C 粘合在一起，然后连接 B、C 的细绳因受扰动而突然断开，弹簧伸展，从而使 C 与 A、B 分离。脱离弹簧后 C 的速度为 $v_{0}$。
\begin{enumerate}
	\item
	求弹簧所释放的势能 $\Delta E$；
	\item
	若更换 B、C 间的弹簧，当物块 A 以初速 $v$ 向 B 运动，物块 C 在脱离弹簧后的速度为 $2v_{0}$，则弹簧所释放的势能 $\Delta E'$ 是多少？
	\item
	若情况（2）中的弹簧与情况（1）中的弹簧相同，为使物块 C 在脱离弹簧后的速度仍为 $2v_{0}$，A 的初速 $v$ 应为多大？
\end{enumerate}
\onepicture[align=right, w=8cm]{\includesvg{figs/28}}
%% answer: 
\jdanswer{
\begin{enumerate}
	\item
	$\Delta E = \frac{1}{3}mv_{0}^{2}$

	\item
	$\Delta E' = \frac{1}{12}m(v - 6v_{0})^{2}$

	\item
	$v = 4v_{0}$，$v = 8v_{0}$（舍去）

\end{enumerate}
}
%% memo:
\memoanswer{
（1）设 A 与 B、C 粘合后的速度为 $u_{1}$。C 脱离弹簧后 A、B 的共同速度为 $u_{2}$。由动量守恒：

$mv_{0} = 3mu_{1}$     ①     $u_{1} = \frac{1}{3}v_{0}$

$2mu_{2} + mv_{0} = mv_{0}$     ②     $u_{2} = 0$

由动能定理：$\frac{1}{2}\cdot 3mu_{1}^{2} + \Delta E = \frac{1}{2}mv_{0}^{2} + \frac{1}{2}\cdot 3mu_{2}^{2}$     ③

$\Delta E = \frac{1}{2}mv_{0}^{2} - \frac{1}{2}\cdot 3m\left(\frac{1}{3}v_{0}\right)^{2} = \frac{1}{3}mv_{0}^{2}$

（2）设与（1）中 $u_{1}$、$u_{2}$ 对应的量为 $u_{1}'$、$u_{2}'$。则有

$mv = 3mu_{1}'$     ④     $u_{1}' = \frac{1}{3}v$

$2mu_{2}' + m\cdot 2v_{0} = mv$     ⑤     $u_{2}' = \frac{1}{2}v - v_{0}$

$\Delta E = \frac{1}{2}m(2v_{0})^{2} + \frac{1}{2}\cdot 2m u_{2}'^{2} - \frac{1}{2}\cdot 3m u_{1}'^{2}$     ⑥

将 $u_{1}'$、$u_{2}'$ 代入⑥式，即得

$\Delta E' = \frac{1}{12}m(v - 6v_{0})^{2}$     ⑦

（3）以 $\Delta E' = \Delta E = \frac{1}{3}mv_{0}^{2}$ 代入⑦式，得 $\frac{1}{3}mv_{0}^{2} = \frac{1}{12}m(v - 6v_{0})^{2}$     ⑧

解得：$v = 4v_{0}$，$v = 8v_{0}$（舍去）。其中 $v = 8v_{0}$ 代入 $u_{2}'$ 表达式得 $u_{2}' = 3v_{0} > 2v_{0}$（C 脱离弹簧后的速度）不合题意，故舍去。

评分标准：全题 15 分。（1）6 分（2）4 分（3）5 分

（1）正确列出①、②、③式各 1 分，正确得出 $u_{1}$、$u_{2}$、$\Delta E$ 各 1 分。

（2）正确列出④式并得出 $u_{1}'$ 共 1 分，正确列出⑤式并得出 $u_{2}'$ 各 1 分，正确得出 $\Delta E'$ 结果 1 分。

（3）正确列出⑧式 1 分，正确得出 $v = 4v_{0}$ 的结果得 2 分，正确说明 $v = 8v_{0}$ 不合理的得 2 分。
}

\end{enumerate}

\end{document}
